\subsection{Clans and additive set functions}

DEFINITION 3. --- A nonempty set \(\Phi\) of subsets of a set \(A\) is said to be a clan if there exists an algebra \(\mathcal{A}\) (over \(\mathbf{R}\) ) consisting of real-valued functions defined on \(A\) , such that the relations \(M \in \Phi\) and \(\varphi_{M} \in \mathcal{A}\) are equivalent.

Example. --- If \(\mu\) is a measure on a locally compact space X then the linear combinations, with real coefficients, of the characteristic functions of integrable sets form an algebra A, because, for any two integrable sets M, N, the function \(\varphi_{M}\varphi_{N} = \varphi_{M\cap N}\) is integrable (No.~5, Prop. 7); it then follows from Defs. 2 and 3 that the set of integrable subsets of X is a clan.

PROPOSITION 17. --- In order that a nonempty set \(\Phi\) of subsets of a set A be a clan, it is necessary and sufficient that it satisfy the following condition:

\begin{enumerate}
\def\labelenumi{(\Roman{enumi})}
\setcounter{enumi}{149}
\tightlist
\item
  For every pair of sets M,N belonging to \(\Phi\) , the sets \(M\cup N\) and \(M\cap CN\) belong to \(\Phi\) .¹
\end{enumerate}

The condition is necessary, by virtue of the relations

\[
\varphi_ {\mathrm{M} \cup \mathrm{N}} = \varphi_ {\mathrm{M}} + \varphi_ {\mathrm{N}} - \varphi_ {\mathrm{M}} \varphi_ {\mathrm{N}}, \quad \varphi_ {\mathrm{M} \cap \mathbf {C N}} = \varphi_ {\mathrm{M}} - \varphi_ {\mathrm{M}} \varphi_ {\mathrm{N}}.
\]

To show that it is sufficient, we first observe that it implies that, for any two sets M,N in \(\Phi\) , M∩N belongs to \(\Phi\) since \(M\cap N=M\cap\mathbb{C}(M\cap\mathbb{C}N)\) . Let \(\mathcal{E}(\Phi)\) be the set of linear combinations, with real coefficients, of the characteristic functions of the sets of \(\Phi\) . Since \(\varphi_{M}\varphi_{N}=\varphi_{M\cap N}\) , \(\mathcal{E}(\Phi)\) is an algebra. Everything comes down to showing that if M is a subset of A such that \(\varphi_{M}=\sum_{i}c_{i}\varphi_{M_{i}}\) , where the \(M_{i}\) belong to \(\Phi\) , then \(M\in\Phi\) . This will result from the following lemma:

Lemma. --- Let \(\Phi\) be a nonempty set of subsets of A satisfying the axiom (CL). Given a finite family \((\mathrm{M}_i)_{1 \leqslant i \leqslant n}\) of sets in \(\Phi\) , there exists a finite family \((\mathbf{N}_j)_{1\leqslant j\leqslant m}\) of pairwise disjoint sets in \(\Phi\) such that each \(\mathbf{M}_i\) is the union of a certain number of the \(\mathbf{N}_j\) .

For, consider the \(2^{n} - 1\) sets of the form \(\bigcap_{i=1}^{n} \mathrm{P}_{i}\) , where \(\mathrm{P}_{i} = \mathrm{M}_{i}\) for certain indices \(i\) , \(\mathrm{P}_{i} = \complement \mathrm{M}_{i}\) for the others, at least one of the \(\mathrm{P}_{i}\) being equal to \(\mathrm{M}_{i}\) . Let \((\mathrm{N}_{j})_{1 \leqslant j \leqslant m}\) be the sequence of these sets arranged in some order; they are pairwise disjoint and belong to \(\Phi\) ; on the other hand, every set \(\mathrm{M}_{k}\) is the union of the sets \(\mathrm{N}_{j} = \bigcap_{i=1}^{n} \mathrm{P}_{i}\) corresponding to the families \((\mathrm{P}_{i})\) such that \(\mathrm{P}_{k} = \mathrm{M}_{k}\) , which proves the lemma.

The lemma established, every function of the form \(\sum_{i=1}^{n} c_i \varphi_{\mathrm{M}_i}\) , where \(\mathrm{M}_i \in \Phi\) , may be written in the form \(\sum_{j=1}^{m} d_j \varphi_{\mathrm{N}_j}\) , where the \(\mathrm{N}_j\) belong to \(\Phi\) and are pairwise disjoint; if \(\varphi_{\mathrm{M}} = \sum_{j=1}^{m} d_j \varphi_{\mathrm{N}_j}\) then necessarily \(d_j = 0\) or \(d_j = 1\) for each index \(j\) , thus \(\mathrm{M}\) is the union of a certain number of the \(\mathrm{N}_j\) , and therefore belongs to \(\Phi\) .

Every clan \(\Phi\) of subsets of A contains the empty subset \(\varnothing\) of A; for, there exists at least one subset \(M \in \Phi\) , therefore \(M - M = \varnothing\) belongs to \(\Phi\) . Note also that the set of subsets of A consisting of the single subset \(\varnothing\) is a clan.

DEFINITION 4. --- Given a clan \(\Phi\) of subsets of a set A, and a Banach space F, one calls step function \(^{2}\) over the sets of \(\Phi\) (or \(\Phi\) -step function), with values in F, every function of the form \(\sum_{i} \mathbf{a}_{i} \varphi_{\mathrm{M}_{i}}\) , where the \(\mathbf{a}_{i}\) belong to F, and the \(\mathrm{M}_{i}\) to \(\Phi\) .

It is clear that the set \(\mathcal{E}_{\mathrm{F}}(\Phi)\) of \(\Phi\) -step functions with values in F is a vector space over R or C. We have just seen in Prop. 17 that the set \(\mathcal{E}(\Phi)\) of real-valued \(\Phi\) -step functions is an algebra over R; it is also the linear subspace of \(R^{A}\) generated by the characteristic functions of the sets of \(\Phi\) .

By the Lemma, every function in \(\mathcal{E}_{\mathrm{F}}(\Phi)\) may be written \(f = \sum_{j} c_{j} \varphi_{N_{j}}\) , where the \(N_{j} \in \Phi\) are pairwise disjoint; from this it follows that \(|f| = \sum_{j} |c_{j}| \varphi_{N_{j}}\) belongs to \(\mathcal{E}(\Phi)\) . In particular, \(\mathcal{E}(\Phi)\) is a Riesz space, since the upper envelope of two functions in \(\mathcal{E}(\Phi)\) belongs to \(\mathcal{E}(\Phi)\) .

Remark. --- It is easily seen that Def. 4 is equivalent to the following: a \(\Phi\) -step function with values in F is a function f that takes on only a finite number of values and which is such that, for every \(a \neq 0\) in F, the set \(\mathbf{f}^{-1}(a)\) belongs to \(\Phi\) .

DEFINITION 5. --- A real-valued function \(\lambda\) defined on a clan \(\Phi\) of subsets of a set \(A\) is said to be additive if, for every pair \(M, N\) of disjoint sets belonging to \(\Phi\) , \(\lambda(M \cup N) = \lambda(M) + \lambda(N)\) .

It follows in particular from this definition that \(\lambda(\varnothing)=0\) .

PROPOSITION 18. --- Let \(\lambda\) be an additive set function defined on a clan \(\Phi\) . There exists one and only one linear form (again denoted \(\lambda\) ) on the vector space \(\mathcal{E}(\Phi)\) of real-valued \(\Phi\) -step functions, such that \(\lambda(\varphi_{\mathrm{M}}) = \lambda(\mathrm{M})\) for every set \(\mathrm{M} \in \Phi\) ; if, moreover, \(\lambda(\mathrm{M}) \geqslant 0\) for every \(\mathrm{M} \in \Phi\) , then \(\lambda\) is a positive linear form on \(\mathcal{E}(\Phi)\) .

The uniqueness of the linear form \(\lambda\) is clear, since the characteristic functions of the sets in \(\Phi\) generate the vector space \(\mathcal{E}(\Phi)\) . To prove the existence of \(\lambda\) , it suffices to prove that the relation \(\sum_{i} c_i \varphi_{\mathrm{M}_i} = 0\) , where the \(\mathrm{M}_i\) are nonempty sets belonging to \(\Phi\) , implies \(\sum_{i} c_i \lambda(\mathrm{M}_i) = 0\) . Now, by the Lemma there exists a finite family \((\mathrm{N}_j)\) of pairwise disjoint nonempty sets in \(\Phi\) such that, for every index \(i\) , \(\varphi_{\mathrm{M}_i} = \sum_{j} a_{ij} \varphi_{\mathrm{N}_j}\) with \(a_{ij} = 0\) or \(a_{ij} = 1\) . The relation \(\sum_{i} c_i \varphi_{\mathrm{M}_i} = 0\) , which may be written \(\sum_{j} \left( \sum_{i} c_i a_{ij} \right) \varphi_{\mathrm{N}_j} = 0\) , therefore implies that \(\sum_{i} c_i a_{ij} = 0\) for every index \(j\) . By Def. 5, we then have

\[
\sum_ {i} c _ {i} \lambda (\mathrm{M} _ {i}) = \sum_ {j} \left(\sum_ {i} c _ {i} a _ {i j}\right) \lambda (\mathrm{N} _ {j}) = 0,
\]

which proves the existence of \(\lambda\) . Finally, suppose that \(\lambda(M) \geqslant 0\) for every \(M \in \Phi\) ; for every function \(f \in \mathcal{E}(\Phi)\) , one can write \(f = \sum_{i} c_i \varphi_{M_i}\) , where the \(M_i \in \Phi\) are pairwise disjoint; if \(f \geqslant 0\) , it follows that \(c_i \geqslant 0\) for every index \(i\) such that \(M_i\) is nonempty, whence \(\lambda(f) = \sum_{i} c_i \lambda(M_i) \geqslant 0\) .

